\documentclass[letterpaper,11pt]{article}

\usepackage[spanish]{babel}
\usepackage[top=2cm,bottom=2cm,right=2cm,left=2cm]{geometry}
\usepackage{amssymb}
\usepackage{amsmath}
\usepackage{amsthm}
\usepackage{empheq}
\usepackage{mdframed}
\usepackage{booktabs}
\usepackage{lipsum}
\usepackage{graphicx}
\usepackage{color}
\usepackage{psfrag}
\usepackage{pgfplots}
\usepackage{bm}
\usepackage{multicol}
\usepackage{wrapfig}
\usepackage{float}
\usepackage{wrapfig}
\usepackage{multirow}
\usepackage{biblatex}
\usepackage{csquotes}
\usepackage{parskip}
\usepackage{cancel}
\usepackage{listing}
\usepackage{tikz}
\usepackage{verbatim}
\usepackage{subcaption}
\usepackage{enumitem}
\usepackage{xcolor}
\usepackage{xfrac}

\graphicspath{{fig/}}
\addbibresource{biblio.bib}
\usepgfplotslibrary{colorbrewer}
\pgfplotsset{width=8cm,compat=1.9}
\definecolor{negro}{RGB}{0,0,0}
\definecolor{azul}{RGB}{0,0,255}
\definecolor{verde}{RGB}{0,255,0}
\definecolor{cian}{RGB}{0,255,255}
\definecolor{rojo}{RGB}{255,0,0}
\definecolor{magenta}{RGB}{255,0,255}
\definecolor{amarillo}{RGB}{255,255,0}
\definecolor{blanco}{RGB}{255,255,255}

\definecolor{limegreen}{cmyk}{0.75,0,1,0}
\definecolor{orange}{cmyk}{0,0.8,0.95,0}
\definecolor{purple}{cmyk}{0.75,1,0,0}

\begin{document}

\onecolumn\begin{@twocolumntrue}
    \begin{minipage}{0.3\textwidth}{\includegraphics[scale=0.24]{Escudo_UD.pdf}} % Cambiar escudo
    \end{minipage}
    \vspace{10pt}
    \begin{minipage}{0.677\textwidth}
        \begin{center}
            \vspace{12mm}
        
            \Large{\textbf{Tarea 3: osciladores y ecuaciones diferenciales}}
            \vspace{3mm}
            
            \large{\textbf{Juan Sebastian Manrique Moreno}}
            \vspace{2mm}
    
            \large{\textit{Ecuaciones diferenciales, Programa Académico de Física, Universidad Distrital Francisco José de Caldas}}
            \vspace{1mm}
            
            Octubre de 2022 % Formato (mes) de (año)
        \end{center}
\vspace{5pt}
\end{minipage}

\centerline{\rule{0.95\textwidth}{0.4pt}} % Línea negra
\end{@twocolumntrue}

\begin{multicols}{2}

\subsection*{\textbf{Transformadas y relaciones}}


Se tienen en cuenta las transformadas y relaciones obtenidas en \cite{R1}, las cuales son:
   
De azimutales a ecuatoriales
\begin{center}
   $(Az, a) \rightarrow  (h, \delta)$
\end{center}
\begin{align}
   \sin h \cos \delta & = \sin Az \cos a \label{eq:sinh} \\ 
   \cos h \cos \delta & =\cos Az \cos a \sin \varphi + \sin a \cos \varphi \label{eq:cosh} \\ 
   \sin \delta & =-\cos Az \cos a \cos \varphi + \sin a \sin \varphi \label{eq:sind}
\end{align}
De ecuatoriales a azimutales:
\begin{center}
   $(h, \delta) \rightarrow (Az, a) $
\end{center}
\begin{align}
   \sin Az \cos a & = \sin h \cos \delta \label{eq:sinaz}\\ 
   \cos Az \cos a & =\cos h \cos \delta \sin \varphi - \sin \delta \cos \varphi \label{eq:cosaz} \\ 
   \sin a & = \cos h \cos \delta \cos \varphi + \sin \delta \sin \varphi \label{eq:sina}
\end{align}
Tiempo sidéreo local:
\begin{align}
   TSL = h_{\star} + \alpha_{\star}
\label{eq:tsl}
\end{align}
   
   %Debido a que las coordenadas tienen un desfase entre el azimut norte ($Az_{N}$) y el azimut sur ($Az_{S}$) se tiene la siguiente relación:
   %\begin{align}
   %Az_{S} = Az_{N} + 180^{\circ}   
   %\end{align}
\section{Eclipse lunar}
   
El próximo\textbf{8 de Noviembre} el máximo del \textbf{eclipse lunar total ocurrirá a las 5:59:11} hora local de Bogotá, cuya latitud es $\varphi = 4^{\circ} ~ 38^{\prime}~ 57.48^{\prime \prime}$. Este eclipse será visible en Pacífico y costa Pacifica de Norte América.
   
\subsection{Inicio del eclipse}
   
El inicio de la totalidad será a las 5:16:39 (averigue el TSL para Bogotá en este instante), si las coordenadas ecuatoriales de la luna para el inicio del eclipse son:
$\alpha = 2~\mathrm{h}~ 48~ \mathrm{min} ~27~ \mathrm{s}, \delta = +16^{\circ}~ 39^{\prime}~ 35^{\prime \prime}$.
   
\textbf{a) Determine las coordenadas azimutales para un observador en Bogotá.}
   
Contando el TSL encontrado en Stellarium, se tienen los siguientes datos:
\begin{align*}
    \alpha &= 2~\mathrm{h}~~ 48~ \mathrm{min}~~ 27~ \mathrm{s} \\
    \delta &= 16^{\circ}~~ 39^{\prime}~~ 35^{\prime \prime}\\
    TSL &= 8~\mathrm{h}~~ 30~\mathrm{min}~~ 40.3~\mathrm{s}
\end{align*}
Se calcula primero el ángulo horario $h$ despejándolo de la ecuación (\ref{eq:tsl}):
\begin{align*}
   h_{\star} &=  TSL - \alpha_{\star} \\
   h_{\star} &=  8.5112~\mathrm{h} - 2.8075~\mathrm{h} \\
   h_{\star} &=  5.7037~\mathrm{h}% ¿aparte la conversión o de una el resultado?
\end{align*}
Se realiza la conversión a grados de $h$:
\begin{align*}
    5.7037~ \cancel{\mathrm{h}} \cdot \frac{15 ^{\circ}}{1~ \cancel{\mathrm{h}}} & =85.5555 ^{\circ} 
\end{align*}
\begin{align}
   h &= 85.5555 ^{\circ} \label{reso:hora}
\end{align}
Se realiza la respectiva conversión para la declinación ($\delta$):
\begin{align*}
   35^{\cancel{\prime \prime}}\cdot \frac{1^{\prime} }{60^{\cancel{\prime \prime}} } &=0.5833^{\prime}\\ 
   39.5833 ^{\cancel{\prime}}\cdot \frac{1 ^{\circ} }{60^{\cancel{\prime}}} &=0.6597^{\circ}
\end{align*}
\begin{align}
    \delta = 16.6597^{\circ} \label{reso:decli}
\end{align}
   
Para hallar \textbf{la altura ($a$)}, se reemplazan los datos en la ecuación(\ref{eq:sina}):
\begin{align*}
    \sin a = &  \cos \left(85.5555^{\circ}\right)  \cos \left(16.6597^{\circ}\right) \cos \left(4.6493^{\circ}\right) + \\
    & \sin \left(16.6597^{\circ}\right) \sin \left(4.6493^{\circ}\right)
\end{align*}
Despejando a de la anterior ecuación se tiene que:
\begin{align*}
   a = 5.58^{\circ}  \Longrightarrow a = 5^{\circ}~~ 34^{\prime}~~ 49^{\prime \prime}
\end{align*}
   
Para el\textbf{ azimuth ($Az$)}, dividiendo las ecuaciones (\ref{eq:sinaz}) y (\ref{eq:cosaz}) resulta que:
\begin{align}
    \tan Az = \frac{\sin h \cos \delta}{\cos h \cos \delta \sin \varphi - \sin \delta \cos \varphi}
\end{align}
   
   % Dividiendo las ecuaciones (\ref{eq:cosaz}) y (\ref{eq:sinaz}) resulta que:
   
   %\begin{align}
       %\cot Az = \frac{\cos h \sen \varphi - \tan %\delta \cos \varphi}{\sin h}
   %\end{align}
   
   Se reemplazan los datos conocidos en la ecuación:
   \begin{align*}
       \tan Az  = & \frac{\sin \left(85.5555^{\circ}\right) \cdot }{\cos\left(85.5555^{\circ}\right)  \cos  \left(16.6597^{\circ}\right)  \sin \left(4.6493^{\circ}\right)}\\
        & \frac{\cos  \left(16.6597^{\circ}\right)}{- \sin \left(16.6597^{\circ}\right) \cos  \left(4.6493^{\circ}\right)}
   \end{align*}
   
   %\begin{align*}
   %\cot Az = & \frac{\cos \left(^{\circ}\right) \sin \left(^{\circ}\right)}{\sin \left(^{\circ}\right)} - \\
   %& \frac{\tan \left(^{\circ}\right) \cos \left(^{\circ}\right)}{\sin \left(^{\circ}\right)}
   %\end{align*}
   
   Despejando $Az$ y realizando una serie de cálculos
   \begin{align}
   Az = -73,677 ^{\circ} \Longrightarrow
   Az = -73 ^{\circ}, 40', 37,2''  %Notación Kartunnen 
   \end{align}
   
   \subsection{Máximo del eclipse}
   
   
   \subsection{Final del eclipse}
   
   El final de la totalidad será a las $6:41:36$ (Averigue el TSL Bogotá), si las coordenadas ecuatoriales de la luna para el fin del
   eclipse son: $\alpha= 2h 51m 31s, \delta = +16^{\circ} 50^{\prime} 53^{\prime \prime}$
   
   \textbf{Coordenadas azimutales para un observador en Bogotá.}
   
   Contando el TSL encontrado en Stellarium, se tienen los siguientes datos:
   \begin{align*}
   \alpha &= 2~\mathrm{h}~~ 51~ \mathrm{min}~~ 31~ \mathrm{s} \\
   \delta &= 16^{\circ}~~ 50^{\prime}~~ 53^{\prime \prime}\\
   TSL &= 7~\mathrm{h}~~ 4~\mathrm{min}~~ 20.3~\mathrm{s}
   \end{align*}
   Se calcula primero el ángulo horario $h$ despejándolo de la ecuación (\ref{eq:tsl}):
   \begin{align*}
       h_{\star} &=  TSL - \alpha_{\star} \\
       h_{\star} &=  9.9309~\mathrm{h} - 2.8586~\mathrm{h} \\
       h_{\star} &=  7.0723~\mathrm{h}% 
   \end{align*}
   Se realiza la conversión a grados de $h$:
   \begin{align*}
   7.0723~ \cancel{\mathrm{h}} \cdot \frac{15 ^{\circ}}{1~ \cancel{\mathrm{h}}} & =106.0845 ^{\circ} 
   \end{align*}
   \begin{align*}
   h &= 106.0845 ^{\circ}
   \end{align*}
   Se realiza la respectiva conversión para la declinación ($\delta$):
   \begin{align*}
   53^{\cancel{\prime \prime}}\cdot \frac{1^{\prime} }{60^{\cancel{\prime \prime}} } &=0.8833^{\prime}\\ 
   50.8833 ^{\cancel{\prime}}\cdot \frac{1 ^{\circ} }{60^{\cancel{\prime}}} &=0.8480^{\circ}
   \end{align*}
   \begin{align*}
   \delta = 16.8480^{\circ}
   \end{align*}
   
   Para hallar \textbf{la altura ($a$)}, se reemplazan los datos en la ecuación(\ref{eq:sina}):
   \begin{align*}
       \sin a = &  \cos \left(106.0845 ^{\circ}\right)  \cos \left(16.8480^{\circ}\right) \cos \left(4.6493^{\circ}\right) + \\
       & \sin \left(16.8480^{\circ}\right) \sin \left(4.6493^{\circ}\right)
   \end{align*}
   Despejando a de la anterior ecuación se tiene que:
   \begin{align*}
   a = -13.934°^{\circ}  \Longrightarrow -13^{\circ}~~ 56^{\prime}~~ 2.4^{\prime \prime}
   \end{align*}
   
   Para el\textbf{ azimuth ($Az$)}, dividiendo las ecuaciones (\ref{eq:sinaz}) y (\ref{eq:cosaz}) resulta que:
   \begin{align}
       \tan Az = \frac{\sin h \cos \delta}{\cos h \cos \delta \sin \varphi - \sin \delta \cos \varphi}
   \end{align}
   
   % Dividiendo las ecuaciones (\ref{eq:cosaz}) y (\ref{eq:sinaz}) resulta que:
   
   %\begin{align}
       %\cot Az = \frac{\cos h \sen \varphi - \tan %\delta \cos \varphi}{\sin h}
   %\end{align}
   
   Se reemplazan los datos conocidos en la ecuación:
   \begin{align*}
       \tan Az  = & \frac{\sin \left(106.0845 ^{\circ}\right) \cdot }{\cos\left(106.0845 ^{\circ}\right)  \cos  \left(16.8480^{\circ}\right)  \sin \left(4.6493^{\circ}\right)}\\
        & \frac{\cos  \left(16.8480^{\circ}\right)}{- \sin \left(16.8480^{\circ}\right) \cos  \left(4.6493^{\circ}\right)}
   \end{align*}
   
   %\begin{align*}
   %\cot Az = & \frac{\cos \left(^{\circ}\right) \sin \left(^{\circ}\right)}{\sin \left(^{\circ}\right)} - \\
   %& \frac{\tan \left(^{\circ}\right) \cos \left(^{\circ}\right)}{\sin \left(^{\circ}\right)}
   %\end{align*}
   
   Despejando $Az$ y realizando una serie de cálculos
   \begin{align}
   Az = -71,35 ^{\circ} \Longrightarrow
   Az = -73 ^{\circ}, 21' %Notación Kartunnen 
   \end{align}
   
   
   
   \section{Coordenadas azimutales y ecuatoriales}
   
   Teniendo en cuenta las coordenadas de los siguientes objetos celestes, realice las operaciones indicadas.
   
   Para observaciones realizadas en Bogotá $(\varphi = 4^{\circ} 38^{\prime} 57.48^{\prime \prime}$ el próximo Equinoccio de Otoño \textbf{22 de septiembre a las 20:03:00}
   
   \textbf{a. Antares:} Es la estrella $\alpha$ Scorpii ($\alpha$ Sco), la estrella más brillante de la constelación de Escorpión, conocida también como
   el corazón del escorpión, es la decimosexta más brillante del cielo nocturno, situada aproximadamente a 550 años luz del
   sistema solar.
   
   i. Sus coordenadas Azimutales para un observador en Bogotá en la fecha y hora indicada serán:
   $Az = 56^{\circ} 46^{\prime} 45^{\prime \prime}, a = 27^{\circ} 40^{\prime} 48^{\prime \prime} $. Hallar las coordenadas ecuatoriales de la estrella.
   \begin{align}
   Az &=^{\circ}~ ^{\prime}~ ^{\prime \prime} \\
   a &= ^{\circ}~ ^{\prime}~ ^{\prime \prime}
   \end{align}
   
   Realizando la respectiva conversión para el azimut ($Az$), resulta:
   \begin{align*}
   ^{\cancel{\prime \prime}}\cdot \frac{1 ^{\prime} }{60  ^{\cancel{\prime \prime}} } &=0.7 '\\ 
    ^{\cancel{\prime}}\cdot \frac{1 ^{\circ} }{60^{\cancel{\prime}}} &=0.645^{\circ}
   \end{align*}
   \begin{align}
   Az = ^{\circ}
   \end{align}
   
   Realizando la respectiva conversión para la altura ($a$), resulta:
   \begin{align*}
   ^{\cancel{\prime \prime}}\cdot \frac{1^{{\prime}} }{60^{\cancel{\prime \prime}}} &= '\\ 
    ^{\cancel{\prime}}\cdot \frac{1 ^{\circ} }{60^{\cancel{\prime}}} &=^{\circ}
   \end{align*}
   \begin{align}
   a = ^{\circ}
   \end{align}
   
   Se halla el ángulo horario dividiendo la ecuación \ref{eq:sinh} entre la ecuación \ref{eq:cosh}
   
   %\begin{align}
   %\cot h = \frac{\cos Az \sen \varphi - \tan a \cos \varphi}{\sin Az}
   %\end{align}
   
   %\begin{align*}
   %\cot h = & \frac{\cos (287,6450 ^{\circ}) \sen (4,6097 ^{\circ})}{\sin (287,6450 ^{\circ})} + \\
   %& \frac{\tan (32,5076  ^{\circ})\cos (4,6097 ^{\circ})}{\sin (287,6450 ^{\circ})}
   %\end{align*}
   
   \textbf{Despejando $h$} y realizando una serie de cálculos
   \begin{align}
   h &= ^{\circ} \\
   h &=  ~ \mathrm{h}~~ ~ \mathrm{min}~~ ~ \mathrm{seg} 
   \end{align}
   
   
   Para hallar  $\delta$ se usa la ecuación \ref{eq:sind}
   
   \begin{align*}
   \sin \delta= & - \cos (^{\circ}) \cos (^{\circ})\\
   & \cos (^{\circ}) + sen(^{\circ})\\ 
   & \sin (^{\circ})
   \end{align*}
   \textbf{Despejando $\delta$ }y realizando una serie de cálculos
   \begin{align}
   \delta =  ^{\circ}  \Longrightarrow
   \delta = ^{\circ}, 
   \end{align}
   
   \textbf{b. Vega de la Lyra:} Es la estrella $\alpha$  Lyra ($\alpha$  Ly), la estrella más brillante de la constelación de la Lyra. Es la quinta estrella más brillante del cielo nocturno . Se considera una estrella relativamente cercana, a solo 25 años luz de la Tierra, siendo una de las más brillantes cercanas al sistema solar.
   
   i. Sus coordenadas Azimutales para un observador en Bogotá en la fecha y hora indicada serán:
   $Az = 150^{\circ} 40^{\prime} 20^{\prime \prime}, a = 49^{\circ} 25^{\prime} 47{\prime \prime} $. Hallar las coordenadas ecuatoriales de la estrella.
   
   \begin{align}
   Az &= 150^{\circ}~ 40^{\prime}~ 20^{\prime \prime} \\
   a &= 49^{\circ}~ 25^{\prime}~ 47^{\prime \prime}
   \end{align}
   
   Realizando la respectiva conversión para el azimut ($Az$), resulta:
   \begin{align*}
   20^{\cancel{\prime \prime}}\cdot \frac{1 ^{\prime} }{60  ^{\cancel{\prime \prime}} } &=0.3333 ^{\prime} \\
   40.3333 ^{\cancel{\prime}}\cdot \frac{1 ^{\circ} }{60^{\cancel{\prime}}} &=0.6722^{\circ}
   \end{align*}
   \begin{align}
   Az = 150,6722 ^{\circ}
   \end{align}
   
   Realizando la respectiva conversión para la altura ($a$), resulta:
   \begin{align*}
   47^{\cancel{\prime \prime}}\cdot \frac{1^{{\prime}} }{60^{\cancel{\prime \prime}}} &= 0.7833 1^{{\prime}\\ 
   25,7833 ^{\cancel{\prime}}\cdot \frac{1 ^{\circ} }{60^{\cancel{\prime}}} &= 0,4297 ^{\circ}
   \end{align*}
   \begin{align}
   a = 49,4297 ^{\circ}
   \end{align}
   
   Se halla el ángulo horario dividiendo la ecuación \ref{eq:sinh} entre la ecuación \ref{eq:cosh}
   
   %\begin{align}
   %\cot h = \frac{\cos Az \sen \varphi - \tan a \cos \varphi}{\sin Az}
   %\end{align}
   
   %\begin{align*}
   %\cot h = & \frac{\cos (287,6450 ^{\circ}) \sen (4,6097 ^{\circ})}{\sin (287,6450 ^{\circ})} + \\
   %& \frac{\tan (32,5076  ^{\circ})\cos (4,6097 ^{\circ})}{\sin (287,6450 ^{\circ})}
   %\end{align*}
   
   \textbf{Despejando $h$} y realizando una serie de cálculos
   \begin{align}
   h &= ^{\circ} \\
   h &=  ~ \mathrm{h}~~ ~ \mathrm{min}~~ ~ \mathrm{seg} 
   \end{align}
   
   Para hallar  $\delta$ se usa la ecuación \ref{eq:sind}
   \begin{align*}
   \sin \delta= & - \cos (^{\circ}) \cos (^{\circ})\\
   & \cos (^{\circ}) + sen(^{\circ})\\ 
   & \sin (^{\circ})
   \end{align*}
   \textbf{Despejando $\delta$ }y realizando una serie de cálculos
   \begin{align}
   \delta =  ^{\circ}  \Longrightarrow
   \delta = ^{\circ}, 
   \end{align}
   
   
   
   
   %ctgAz = ctgh senphi - tandelta (cosphi/senh)
   %=> ctgAz = (cosh senphi - tandelta cosphi)/senh
   %ctgh = ctgAz senphi + tana (cosphi/senAz)
   
   \selectlanguage{spanish}
   \begin{thebibliography}{99}
   
   
   \bibitem{R1} 
   
   Manrique, J. S. Triana, L. C. (2022). Quiz 02. Astronomía y Laboratorio. Universidad Distrital FJDC.
   
   
   \end{thebibliography}
   
   \end{multicols}
   
   \end{document}